\scanpage{11}
Il faut prouver qu'on a bien un foncteur.

a) $\varphi_{X,X}(\mathrm{id}_X)=\Delta_X$. En effet, \deleted{\unsure{le produit est}} \inserted{c'est la \uline{définition}} de $\Delta_X$ !

\diagram{\begin{tikzcd}[ampersand replacement=\&]
\text{b)}\quad X\arrow[r,dash,"f"]\&Y\arrow[r,dash,"g"]\&Z\qquad\varphi(g)\varphi(f)=\varphi(gf)
\end{tikzcd}}

Le premier mb est

\diagram{\begin{tikzcd}[ampersand replacement=\&,column sep=small,row sep=large]
\&X\times Z\arrow[dl,"\mu'"']\arrow[dr,"\lambda'"]\&\&X\times Y\arrow[dl,"{\unsure{\lambda^{?}|\lambda}=\mathrm{id}_X\times\unsure{\Gamma|\mathscr T}_g}"']\arrow[dr,"\mu"]\&\\
X\arrow[dr,"\Gamma_f"']\&\&X\times Y\times Z\arrow[dl,"p_{21}"']\arrow[d,"p_{31}"]\arrow[dr,"p_{32}"]\&\&Y\arrow[dl,"\Gamma_g"]\\
\&X\times Y\&X\times Z\&Y\times Z\&
\end{tikzcd}}

\[
p_{31*}\Bigl(
\underbrace{p_{32}^*(\varphi(g))}_{\substack{\Vert\\p_{32}^*(\Gamma_{g*}(1))\\\Vert\\\deleted{\unsure{\mu_*|p_*}}\lambda_*(\mu^*(1))=\lambda_*(1)}}
\underbrace{p_{21}^*(\varphi(f))}_{\substack{\Vert\\p_{21}^*\Gamma_{f*}(1)\\\Vert\\\lambda'_*(\mu'{}^*(1))=\lambda'_*(1)}}
\]

\[
\varphi(g)\varphi(f)=p_{31*}\bigl(\lambda^{\unsure{+|*|{}}}_*(1)\,\lambda'_*(1)\bigr)\qquad\text{Mais}
\]

\diagram{\begin{tikzcd}[ampersand replacement=\&,column sep=small,row sep=large]
\&\&X\arrow[dl,hook,"\Gamma_{gf}"']\arrow[dr,hook,"\Gamma_f"]\&\&\\
\&X\times Z\arrow[dl,"\mu'=p_X"']\arrow[dr,hook,"{\lambda'=\deleted{\unsure{I}}\Gamma_f\times Z}"']\&\text{(cartésien)}\&X\times Y\arrow[dl,hook,"{\lambda=\mathrm{id}_X\times\Gamma_g}"]\arrow[dr,"\mu=p_Y"]\&\\
X\arrow[dr,hook,"\Gamma_f"']\&\&X\times Y\times Z\arrow[dl,"p_{21}"']\arrow[d,"p_{31}"]\arrow[dr,dash,"p_{32}"]\&\&Y\arrow[dl,hook,"\Gamma_g"]\\
\&X\times Y\&X\times Z\&Y\times Z\&
\end{tikzcd}}

\[
\begin{aligned}
\lambda_*(1)\,\unsure{{}_*|{}}\lambda'_*(1)&=\lambda_*\bigl(\underbrace{\lambda^*(\lambda'_*(1))}_{\substack{\Vert\\\Gamma_{f*}(\Gamma_{gf}^*(1))=\Gamma_{f*}(1)}}\bigr)\quad\boxed{\text{form. proj}}\\
&=\lambda_*\Gamma_{f*}(1)=\unsure{\nu|v}_*(\unsure{1|1/1})
\end{aligned}
\]

où $\nu:X\mathbin{-}X\times Y\times Z$ est donné par $\nu(x)=(x,f(x),gf(x))$

\[
\begin{aligned}
x&\mapsto (x,gf(x))\unsure{\vdash|\mapsto|\leftarrow} (x,f(x),gf(x))\\
x&\vdash (x,f(x))\vdash (x,f(x),\unsure{gf(x))}
\end{aligned}
\]

\[
\begin{aligned}
&\deleted{\unsure{x\mapsto{}}}\ (x,f(x),z)\\
&(x,y,g(y))\\
&(x,f(x),gf(x))
\end{aligned}
\]

D'où

\[
\varphi(g)\varphi(f)=p_{31*}\bigl(\nu_*(1)\bigr)=\Gamma_{gf*}(1)=\varphi(gf)\qquad\text{OK.}
\]

\noindent\rule{3.5cm}{0.3pt}\par\nobreak\vspace{-0.8\baselineskip}\noindent\rule{3.6cm}{0.3pt}

Définissons le foncteur $\psi$, \inserted{avec $\psi(X)=A(\unsure{X|x})$,} i.e. pour $u\in\operatorname{Hom}_{\unsure{V|\mathcal V}^{\mathcal A}}(X,Y)=A(X\times Y)$, et $x\in A(X)$, définissons $u(x)\in A(Y)$ par \unsure{produit|la formule} \deleted{\unsure{$u(x)$}}

\[
u(x)=p_{2*}(u\,p_1^*(x))
\]

Prouvons que l'on a bien une \uline{loi fonctorielle}.

% archive sheet number: 10
